\documentclass[11pt,twocolumn]{article}

\usepackage[margin=1in]{geometry}
\usepackage{amsmath,amssymb,amsfonts}
\usepackage{booktabs}
\usepackage{hyperref}
\usepackage{listings}
\usepackage{xcolor}
\usepackage{enumitem}
\usepackage{caption}

\hypersetup{colorlinks=true,linkcolor=blue,citecolor=blue,urlcolor=blue}

\lstset{
  basicstyle=\ttfamily\small,
  keywordstyle=\color{blue}\bfseries,
  commentstyle=\color{gray},
  breaklines=true,
  frame=single,
  columns=fullflexible
}

\title{\textbf{LLVM2: Unified Solver-Driven Compilation\\Across Heterogeneous Compute Targets}\\[0.5em]
\large A Verified Compiler Backend with Proof-Carrying IR}

\author{Andrew Yates\\Dropbox, Inc.\\{\small\texttt{ayates@dropbox.com}}}

\date{April 2026}

\begin{document}
\maketitle

\begin{abstract}
We present LLVM2, a verified compiler backend that uses SMT solvers to simultaneously optimize and verify code generation across heterogeneous compute targets---CPU, SIMD, GPU, and neural accelerators. Unlike traditional compilers that rely on thousands of hand-written pattern-matching rules, LLVM2 uses a single solver-driven architecture: the solver \emph{synthesizes} optimal instruction sequences across all targets and \emph{proves} each transformation correct. The key enabler is proof-carrying intermediate representation (tMIR), which provides formal annotations---purity, non-aliasing, bounds safety, associativity---that eliminate the conservative assumptions forcing traditional compilers to keep computations on the CPU. We show that superoptimization, heterogeneous compute allocation, compilation transparency, and AI-native extensibility are not separate features but manifestations of a single unified architecture: one solver, all targets, every decision proven and justified.
\end{abstract}

\section{Introduction}

Modern compilation is built on two independent traditions: \emph{optimization} (pattern-matching rewrite rules applied by hand-written passes) and \emph{verification} (proving compiler correctness after the fact, as in CompCert~\cite{compcert} and Alive2~\cite{alive2}). Meanwhile, \emph{heterogeneous compute allocation}---deciding whether a computation runs on CPU, GPU, or accelerator---is left entirely to the programmer via explicit APIs (Metal, CUDA, CoreML).

These three concerns---optimization, verification, and resource allocation---are traditionally treated as separate problems. We argue they are the \emph{same problem} and present a unified architecture that solves all three simultaneously using an SMT solver.

\subsection{The Core Insight}

An SMT solver does not care what the compilation target is. Given semantic encodings of a source computation and a candidate implementation, it checks equivalence:
\begin{equation}
\forall\, \mathit{inputs}: \mathit{source}(\mathit{inputs}) = \mathit{candidate}(\mathit{inputs})
\label{eq:equiv}
\end{equation}

If the candidate is an AArch64 scalar sequence, the solver verifies scalar code. If the candidate is a NEON SIMD sequence, a Metal GPU kernel, or a Neural Engine operation, the solver verifies that too. The solver searches over all targets simultaneously and selects the cheapest proven-correct implementation.

The missing ingredient that makes this tractable is \emph{proof-carrying IR}. Without formal annotations, a compiler cannot safely move computations to GPU because it must conservatively assume aliasing, out-of-bounds access, and non-determinism. With proof annotations from tMIR---purity, non-aliasing (\textsc{ValidBorrow}), bounds safety (\textsc{InBounds}), associativity---the solver can prove cross-target equivalence for computations that would be illegal to move in a traditional compiler.

\subsection{Contributions}

\begin{enumerate}[leftmargin=*,topsep=4pt,itemsep=2pt]
\item A \textbf{unified solver architecture} where superoptimization, heterogeneous compute allocation, transparency, and AI extensibility are aspects of a single system (Section~\ref{sec:architecture}).
\item A \textbf{proof-guided target analysis} showing how tMIR proof annotations prune the cross-target search space from intractable to feasible (Section~\ref{sec:proofs}).
\item \textbf{Multi-target semantic encoding} extending SMT-based verification from scalar instructions to SIMD, GPU kernels, and neural accelerator operations (Section~\ref{sec:encoding}).
\item A \textbf{self-improving compilation loop} where AI agents propose optimizations, the solver verifies them, and proven-correct rules are added automatically (Section~\ref{sec:ai}).
\end{enumerate}

\section{Background and Related Work}

\subsection{Superoptimization}

STOKE~\cite{stoke} uses Markov Chain Monte Carlo to stochastically search for optimal x86-64 instruction sequences, with three-tier verification: testing, bounded checking, and full SMT proof. Starting from \texttt{clang -O0} output, STOKE produces code matching \texttt{gcc -O3}.

Souper~\cite{souper} uses enumerative synthesis with Counter-Example Guided Inductive Synthesis (CEGIS) to find peephole optimizations in LLVM IR. It achieved a 4.4\% reduction in Clang binary size by discovering optimizations that LLVM's hand-written rules missed.

Both systems target a single ISA and verify post-hoc. LLVM2 integrates verification into the search loop and searches across multiple targets simultaneously.

\subsection{Verified Compilation}

CompCert~\cite{compcert} provides end-to-end Coq proofs for a C compiler. Alive2~\cite{alive2} uses Z3 to verify individual LLVM IR transformations, finding 40+ miscompilation bugs.

LLVM2 combines synthesis and verification: instead of verifying human-written rules, the solver \emph{discovers} rules and proves them correct by construction.

\subsection{Heterogeneous Compilation}

Halide~\cite{halide} separates algorithm from schedule, enabling automated search over execution strategies including CPU, GPU, and vector targets. TVM~\cite{tvm} and Ansor~\cite{ansor} use learned cost models for tensor operation autotuning.

These systems are domain-specific (image processing, ML inference). LLVM2 targets general-purpose computation with formal correctness guarantees.

\subsection{AI-Guided Compilation}

MLGO~\cite{mlgo} replaces LLVM heuristics with neural networks for inlining and register allocation (3--7\% binary size reduction). CompilerGym~\cite{compilergym} provides an RL environment for optimization pass ordering.

LLVM2 goes further: AI agents can propose \emph{arbitrary} optimization rules, and the solver serves as a formal gatekeeper---only proven-correct rules are accepted.

\section{Unified Solver Architecture}
\label{sec:architecture}

\subsection{Overview}

LLVM2 compiles from tMIR (proof-carrying IR) to machine code through five phases, all mediated by the SMT solver z4:

\begin{enumerate}[leftmargin=*,topsep=4pt,itemsep=2pt]
\item \textbf{Proof-guided target analysis.} For each computation subgraph, examine tMIR proof annotations to determine which compute targets (CPU, SIMD, GPU, ANE) are legal.
\item \textbf{Solver-driven synthesis.} For each subgraph and legal target, enumerate candidate implementations, prove equivalence via CEGIS, rank by cost model.
\item \textbf{Dispatch plan generation.} Partition the computation across selected targets, generate data transfer and synchronization code, prove the plan preserves semantics.
\item \textbf{Per-target lowering.} Lower each partition to its target: AArch64 for CPU, NEON for SIMD, Metal IR for GPU, CoreML for ANE. Each lowering verified by z4.
\item \textbf{Transparency and certification.} Emit proof certificates, event logs, and provenance maps as compilation artifacts.
\end{enumerate}

The key architectural property is that every phase uses the \emph{same} solver and the \emph{same} equivalence query (Equation~\ref{eq:equiv}). There is no separate ``optimization engine'' and ``verification engine''---they are the same system.

\subsection{The Resource Allocation Analogy}

Register allocation maps virtual registers to physical registers (X0--X30), minimizing spills to memory. This is NP-hard but solved routinely by compilers.

Compute allocation maps computation subgraphs to compute targets (CPU, SIMD, GPU, ANE), minimizing total execution time including data transfer overhead. This is graph partitioning---the same complexity class.

\begin{table}[h]
\centering
\caption{Register allocation vs.\ compute allocation}
\label{tab:analogy}
\small
\begin{tabular}{@{}lll@{}}
\toprule
\textbf{Concept} & \textbf{Register Alloc.} & \textbf{Compute Alloc.} \\
\midrule
Resources & Registers (X0--X30) & CPU, GPU, ANE, SIMD \\
Values & Virtual registers & Computation subgraphs \\
Constraints & Live ranges, ABI & Transfer cost, launch overhead \\
Spilling & Register to stack & GPU to CPU fallback \\
Cost model & Instruction latency & Compute + transfer time \\
\bottomrule
\end{tabular}
\end{table}

The compiler already solves the hard version (register allocation). Compute allocation is the same structure at coarser granularity. The programmer should not manually assign either.

\section{Proof-Guided Target Analysis}
\label{sec:proofs}

The cross-target search space is intractable without pruning. tMIR proof annotations prune aggressively:

\begin{table}[h]
\centering
\caption{tMIR proof annotations and their compilation effects}
\label{tab:proofs}
\small
\begin{tabular}{@{}lll@{}}
\toprule
\textbf{Annotation} & \textbf{Eliminates} & \textbf{Enables} \\
\midrule
\textsc{Pure} & ``May have side effects'' & Move to any target \\
\textsc{ValidBorrow} & ``Pointers may alias'' & Zero-copy DMA \\
\textsc{InBounds} & ``Index may be OOB'' & Unchecked GPU indexing \\
\textsc{NoOverflow} & ``May overflow'' & Fast GPU arithmetic \\
\textsc{Commutative} & ``Order matters'' & GPU-friendly reorder \\
\textsc{Associative} & ``Grouping matters'' & Parallel reductions \\
\bottomrule
\end{tabular}
\end{table}

\subsection{Pruning Algorithm}

Given a computation subgraph $G$:

\begin{enumerate}[leftmargin=*,topsep=4pt,itemsep=2pt]
\item If $G$ is not \textsc{Pure}: restrict to CPU only.
\item If $G$ does not operate on arrays: restrict to scalar/SIMD.
\item If array accesses lack \textsc{InBounds} and \textsc{ValidBorrow}: restrict to CPU/SIMD.
\item If $G$ is \textsc{Associative} and \textsc{Commutative}: enable parallel GPU reductions.
\item Apply cost model: if data volume is below GPU launch threshold, prefer SIMD.
\end{enumerate}

Steps 1--3 use only proof annotations (zero solver cost). The solver runs only on candidates surviving this filter---typically less than 5\% of the full search space.

\subsection{Why Traditional Compilers Cannot Do This}

Without proof annotations, a compiler must assume the worst case at each step. Without \textsc{Pure}, the computation might write to global state, so moving it to GPU changes observable behavior. Without \textsc{ValidBorrow}, arrays might alias, corrupting data during transfer. Without \textsc{InBounds}, GPU indexing might access out-of-bounds memory differently than CPU. Without \textsc{Associative}, parallel reduction on GPU might produce a different result than sequential CPU computation.

This is why GPU programming requires explicit programmer annotation (Metal shaders, CUDA kernels, CoreML models). The compiler lacks the information to make these decisions safely. tMIR provides that information.

\section{Multi-Target Semantic Encoding}
\label{sec:encoding}

\subsection{AArch64 Scalar (Existing)}

Each instruction is encoded as a bitvector function:
\begin{align*}
[\![ \mathtt{ADD}\ X_d, X_n, X_m ]\!] &= \lambda(X_n, X_m).\ \mathit{BVAdd}(X_n, X_m) \\
[\![ \mathtt{LSL}\ X_d, X_n, \#k ]\!] &= \lambda(X_n).\ \mathit{BVShl}(X_n, k)
\end{align*}

Theory: QF\_BV (quantifier-free bitvectors). This is standard and well-supported.

\subsection{NEON SIMD (New)}

128-bit vector operations decomposed into lanes:
\begin{align*}
[\![ \mathtt{FADD.2D}\ &V_d, V_n, V_m ]\!] = \lambda(V_n, V_m).\\
&\mathit{concat}(\mathit{FPAdd}(V_n[63{:}0], V_m[63{:}0]),\\
&\qquad\quad\ \mathit{FPAdd}(V_n[127{:}64], V_m[127{:}64]))
\end{align*}

Theory: QF\_BV + QF\_FP (floating-point).

\subsection{GPU Kernel (New)}

Parallel map and reduce over arrays:
\begin{align*}
[\![ \mathit{pmap}(f, \mathit{arr}) ]\!] &= \lambda(\mathit{arr}).\ \forall i \in [0, n): \mathit{result}[i] = f(\mathit{arr}[i]) \\
[\![ \mathit{preduce}(\oplus, \mathit{arr}) ]\!] &= \lambda(\mathit{arr}).\ \mathit{fold}(\oplus, 0, \mathit{arr})
\end{align*}

The parallel reduce encoding requires \textsc{Associative} and \textsc{Commutative} proofs from tMIR to justify the equivalence with sequential evaluation.

Theory: QF\_ABV (arrays + bitvectors) with bounded quantification.

\subsection{Neural Engine (New)}

Matrix operations with fixed-function semantics:
\[
[\![ \mathit{matmul}(A, B) ]\!] = \lambda(A, B).\ \forall i,j: R[i][j] = \textstyle\sum_k A[i][k] \cdot B[k][j]
\]

Theory: Nested array theory or unrolled bitvectors for bounded dimensions.

\section{Solver-Driven Synthesis}

\subsection{CEGIS Loop}

For each candidate implementation $c$ of source computation $s$:

\begin{enumerate}[leftmargin=*,topsep=4pt,itemsep=2pt]
\item Encode $s$ and $c$ as SMT formulas.
\item Query z4: does there exist an input where $s$ and $c$ disagree?
\item If UNSAT: $c$ is proven equivalent. Record proof. Rank by cost.
\item If SAT: z4 returns counterexample. Add to test suite. Try next candidate.
\end{enumerate}

\subsection{Search Strategy}

We combine enumerative synthesis (Souper) for systematic coverage with stochastic search (STOKE) for escaping local optima:

\begin{itemize}[leftmargin=*,topsep=4pt,itemsep=2pt]
\item \textbf{Offline:} Enumerate all sequences up to length 3 across all targets. Run at build time. Cache proven rules.
\item \textbf{Online:} For hot paths, time-budgeted stochastic search (100ms per basic block). Apply STOKE-style mutations across targets.
\end{itemize}

\subsection{Cross-Target Synthesis}

The search space includes hybrid implementations: CPU scalar + NEON vectorized inner loop; CPU dispatch + GPU kernel + CPU reduction; CPU pre-processing + ANE matrix multiply + CPU post-processing.

The solver verifies the \emph{composition}: the data transfer, kernel execution, and synchronization together preserve the source semantics.

\subsection{Cost Model}

Candidates are ranked by estimated execution time:
\[
\mathit{cost}(c) = \mathit{compute}(c) + \mathit{transfer}(c) + \mathit{launch}(c)
\]

Phase 1 uses hand-written tables (Apple M-series data from microarchitecture studies~\cite{m1}). Phase 2 uses learned cost models following Ithemal~\cite{ithemal}.

\section{AI-Native Self-Improving Compilation}
\label{sec:ai}

AI agents (RL, LLM, or evolutionary search) propose optimization rules as (pattern, replacement) pairs. The solver verifies each proposal. If UNSAT (no distinguishing input exists), the rule is proven correct and added to the database. If SAT, the counterexample is returned to the agent.

The solver is the \emph{formal gatekeeper}: AI agents can be arbitrarily creative or wrong, but only proven-correct rules enter the database. The system can only improve, never regress in correctness.

The self-improving loop: compile programs, measure runtime, AI agent proposes new rules from observing IR and cost feedback, solver verifies, proven rules are added, performance is re-measured. A compiler that continuously discovers new optimizations with formal correctness guarantees and zero human intervention.

\section{Transparency as Inherent Property}

In a solver-driven architecture, transparency is not a feature to be added---it is an inherent property. The solver produces proofs for every equivalence check (proof certificates). The cost model produces numeric justifications for every target selection (cost justification). Each solver query links source tMIR to target machine instruction (provenance). Every candidate considered and every rule applied or rejected is recorded (decision log).

This contrasts with traditional compilers where transparency (e.g., LLVM's opt-remarks~\cite{llvm-remarks}) is bolted on and provides limited ``what'' without ``why.''

\section{Implementation Status}

LLVM2 is implemented in Rust (61K lines, 1,688 tests) across six crates. Currently operational: AArch64 scalar compilation (50 opcodes), 174 SMT proofs, Mach-O object emission. Toy integer programs compile and execute correctly on Apple Silicon.

The unified solver architecture, multi-target encoding, and AI-native loop are designed but not yet implemented. The critical dependency is z4 array theory (QF\_ABV) for GPU semantic encoding.

\section{Conclusion}

We presented LLVM2, a compiler backend built on the insight that optimization, verification, and heterogeneous compute allocation are the same problem. A single SMT solver searches across all targets, proves every transformation correct, and produces auditable certificates. Proof-carrying IR is the key enabler---it provides the formal annotations that eliminate the conservative assumptions trapping traditional compilers on the CPU.

The programmer writes pure math. The compiler maps it to the best hardware. With proofs.

\begin{thebibliography}{99}
\small

\bibitem{stoke}
E.~Schkufza, R.~Sharma, and A.~Aiken.
Stochastic superoptimization.
In \emph{ASPLOS}, 2013.

\bibitem{souper}
R.~Sasnauskas et al.
Souper: A synthesizing superoptimizer.
\emph{arXiv:1711.04422}, 2017.

\bibitem{compcert}
X.~Leroy.
Formal verification of a realistic compiler.
\emph{Comm.\ ACM}, 52(7):107--115, 2009.

\bibitem{alive2}
N.~P.~Lopes, J.~Lee, C.-K.~Hur, Z.~Liu, and J.~Regehr.
Alive2: Bounded translation validation for LLVM.
In \emph{PLDI}, 2021.

\bibitem{halide}
J.~Ragan-Kelley et al.
Halide: Optimizing parallelism, locality, and recomputation.
In \emph{PLDI}, 2013.

\bibitem{tvm}
T.~Chen et al.
TVM: An automated end-to-end optimizing compiler for deep learning.
In \emph{OSDI}, 2018.

\bibitem{ansor}
L.~Zheng et al.
Ansor: Generating high-performance tensor programs.
In \emph{OSDI}, 2020.

\bibitem{mlgo}
M.~Trofin et al.
MLGO: A machine learning guided compiler optimizations framework.
\emph{arXiv:2101.04808}, 2022.

\bibitem{compilergym}
C.~Cummins et al.
CompilerGym: Robust compiler optimization environments for AI research.
In \emph{CGO}, 2022.

\bibitem{ithemal}
C.~Mendis, A.~Renda, S.~Amarasinghe, and M.~Carbin.
Ithemal: Accurate basic block throughput estimation using deep neural networks.
In \emph{ICML}, 2019.

\bibitem{m1}
D.~Johnson.
Apple M1 microarchitecture research.
\emph{dougallj.github.io/applecpu/firestorm.html}, 2021.

\bibitem{llvm-remarks}
LLVM Project.
LLVM optimization remarks.
\emph{llvm.org/docs/Remarks.html}, 2024.

\bibitem{egg}
M.~Willsey et al.
egg: Fast and extensible equality saturation.
In \emph{POPL}, 2021.

\end{thebibliography}

\end{document}
